\documentclass[11pt]{article}
\usepackage[T1]{fontenc}
\usepackage{lmodern,amsmath,amssymb,amsthm,booktabs,microtype}
\usepackage[margin=1in]{geometry}
\usepackage[hidelinks]{hyperref}
\newtheorem{theorem}{Theorem}[section]
\newtheorem{lemma}[theorem]{Lemma}
\newtheorem{proposition}[theorem]{Proposition}
\newcommand{\F}{\mathbb F_2}
\newcommand{\OG}{\mathcal O}
\newcommand{\Span}{\operatorname{span}}
\title{Clique and chromatic numbers of binary subspace orthogonality graphs}
\author{Haobo Ma \and Reza Nikandish \and Wenlin Zhang}
\date{Working draft, 29 September 2026}
\begin{document}
\maketitle
\begin{abstract}
We determine the clique number of the orthogonality graph on the nonzero
subspaces of $\F^n$ and establish that its chromatic number in dimension six
is $15$. The clique upper bound follows from a radical and quotient argument.
The chromatic upper bound is certified by a complete coloring of all $2\,824$
nonzero subspaces and an independent check against coordinate orthogonality.
For the induced graph on lines, we prove a weighted independent-set bound and
construct a $13$-coloring, giving $12\leq\chi(\Gamma_6)\leq13$.
Thus the full subspace graph requires at least two more colors than its line
subgraph. We describe the geometry of the fixed-clique color lists and a
dimension reduction for extending the search to dimension seven.
\end{abstract}
\section{Introduction and statements of the results}
\label{sec:intro}

Let \(V_n=\F^n\) be equipped with the standard bilinear form
\[
\langle x,y\rangle=\sum_{i=1}^n x_i y_i .
\]
The \emph{orthogonality graph} \(\OG_n^*\) has as its vertices the
nonzero linear subspaces of \(V_n\).  Two distinct vertices \(U,W\)
are adjacent when every vector of \(U\) is orthogonal to every vector
of \(W\).  Its induced subgraph on the one-dimensional subspaces is
denoted by \(\Gamma_n\).  Over \(\F\), each line has a unique nonzero
vector, so \(\Gamma_n\) can equivalently be described as the graph on
the nonzero vectors of \(V_n\) with two distinct vectors adjacent when
their dot product is zero.

The graph \(\OG_n^*\) was introduced by Nikandish \cite{Nikandish2026}
in connection with the annihilating-ideal graph of the ring
\[
R_n=\frac{\F[x_1,\dots,x_n,z]}{(z^2,zx_i,x_ix_j\,(i\ne j),x_i^2-z)} .
\]
Indeed, the annihilating-ideal graph \(\AG(R_n)\) is isomorphic to
\(\OG_n\).  Nikandish \cite{Nikandish2026} posed the following
problems.

\begin{problem}\label{prob:clique}
Determine the exact value of \(\omega(\OG_n^*)\) for every \(n\).
\end{problem}

\begin{problem}\label{prob:chromatic}
Determine \(\chi(\OG_n^*)\) for \(n\ge6\) and investigate the relation
between \(\chi(\OG_n^*)\) and \(\omega(\OG_n^*)\).
\end{problem}

\begin{problem}\label{prob:infinite}
Determine whether there exist infinitely many \(n\ge1\) with
\(\chi(\OG_n^*)>\omega(\OG_n^*)\).
\end{problem}

Put
\[
N(r)=\sum_{d=1}^{r} {r\brack d}_2 ,
\]
where \({r\brack d}_2\) is the Gaussian binomial coefficient, and set
\(N(0)=0\).  Our results are
\begin{equation}\label{eq:clique-formula}
\omega(\OG_n^*)=\max\bigl\{n,\ N(\lfloor n/2\rfloor)+(n\bmod 2)\bigr\}
\qquad(n\ge1),
\end{equation}
and
\begin{equation}\label{eq:chromatic-six}
\chi(\OG_6^*)=15,\qquad 12\le\chi(\Gamma_6)\le13 .
\end{equation}
The first formula answers Problem~\ref{prob:clique} in all
dimensions.  The second answers the dimension-six case of
Problem~\ref{prob:chromatic} and quantifies the separation between
coloring the line graph and coloring the full subspace graph.

The clique argument is algebraic: it sums the radicals of an arbitrary
clique and passes to the resulting orthogonal quotient.  The
full-graph chromatic upper bound uses an explicit finite certificate,
while the line lower bound follows from the geometry of pairwise
nonorthogonal vectors.  The verification details and the scope of the
Lean formalization are collected in Section~\ref{sec:computational}.

The paper is organized as follows.  Section~\ref{sec:clique} contains
the proof of \eqref{eq:clique-formula}.
Section~\ref{sec:line} treats the line graph \(\Gamma_6\) and
establishes the bounds in \eqref{eq:chromatic-six}.
Section~\ref{sec:chromatic} proves \(\chi(\OG_6^*)=15\).
Section~\ref{sec:geometry} describes the geometry of the precoloring.
Section~\ref{sec:higher} records the dimension reduction for \(n=7\)
and states the remaining open problems.
Section~\ref{sec:computational} collects the computational
certificates and the Lean formalization.

\section{The exact clique number}
\label{sec:clique}

Let \(N(r)\) denote the number of nonzero subspaces of \(\F^r\), with
\(N(0)=0\).  Equivalently, \(N(r)=\sum_{d=1}^r {r\brack d}_2\).

\begin{theorem}\label{thm:clique}
For every positive integer \(n\),
\[
\omega(\OG_n^*)=\max\bigl\{n,\ N(\lfloor n/2\rfloor)+(n\bmod 2)\bigr\} .
\]
\end{theorem}

\begin{proof}
We prove the upper and lower bounds separately.

\emph{Lower bound.}  The \(n\) coordinate lines
\(\langle e_1\rangle,\dots,\langle e_n\rangle\) are pairwise
orthogonal and form a clique of size \(n\).  For the second term, let
\(r=\lfloor n/2\rfloor\) and take
\[
T=\Span\{e_1+e_2,\ e_3+e_4,\ \dots,\ e_{2r-1}+e_{2r}\} .
\]
The form \(\langle\cdot,\cdot\rangle\) vanishes on \(T\), so every
nonzero subspace of \(T\) is orthogonal to every other one.  Hence the
nonzero subspaces of \(T\) form a clique of size \(N(r)\).  When \(n\)
is odd, \(\dim T^\perp=r+1\), so \(T^\perp\) is a distinct additional
vertex orthogonal to every subspace of \(T\).  These two
constructions attain both terms of the maximum.

\emph{Upper bound.}  Let \(\mathcal C\) be a clique in \(\OG_n^*\).
For each \(U\in\mathcal C\) put
\[
\rad(U)=U\cap U^\perp ,
\]
and define
\[
R=\sum_{U\in\mathcal C}\rad(U),\qquad r=\dim R .
\]
By construction, \(R\subseteq U^\perp\) for every \(U\in\mathcal C\),
and the summands are pairwise orthogonal.  Hence \(R\subseteq R^\perp\),
so \(R\) is totally isotropic and \(2r\le n\).  Moreover,
\[
U\subseteq R^\perp,\qquad R\subseteq R^\perp,\qquad
U\cap R=\rad(U)\qquad(U\in\mathcal C).
\]
For the last equality, one inclusion is the definition of \(R\); the
other follows from \(R\subseteq U^\perp\).

Since \(B\) is nondegenerate, \(\dim(R^\perp/R)=n-2r\).  The form
descends to the quotient \(Q=R^\perp/R\).  Write \(\overline U\) for
the image of \(U\) in \(Q\).  We claim that the restricted form on
\(\overline U\) is nondegenerate.  Indeed, if \(u\in U\) represents a
vector annihilating \(\overline U\), then \(u\in\rad(U)\subseteq R\),
so its image in \(Q\) is zero.

The spaces \(\overline U\) are pairwise orthogonal in \(Q\).  Their
nonzero members form an indexed direct sum: in any relation
\(\sum_U x_U=0\) with \(x_U\in\overline U\), pairing with arbitrary
vectors of one \(\overline W\) and using the nondegeneracy of the
restriction to \(\overline W\) gives \(x_W=0\).  Hence there are at
most \(\dim Q=n-2r\) nonzero images.

The members of \(\mathcal C\) with zero image are precisely those
contained in \(R\), and there are at most \(N(r)\) of them.  Therefore
\begin{equation}\label{eq:bound}
|\mathcal C|\le N(r)+n-2r .
\end{equation}
For \(r=0\), this gives \(|\mathcal C|\le n\).  For every \(s\ge1\),
embedding \(\F^s\) as a hyperplane in \(\F^{s+1}\) gives
\[
N(s+1)\ge N(s)+2 ,
\]
because an outside line and the whole ambient space are two distinct
additional subspaces.  Hence \(N(s)-2s\) is nondecreasing for
\(s\ge1\).  If \(r\ge1\), putting \(t=\lfloor n/2\rfloor\) in
\eqref{eq:bound} gives
\[
|\mathcal C|\le N(t)+n-2t=N(t)+(n\bmod 2).
\]
This completes the proof of the upper bound and hence of the
theorem.
\end{proof}

\begin{corollary}\label{cor:15clique}
For \(n=6\),
\[
\omega(\OG_6^*)=N(3)=15 .
\]
An explicit clique of size \(15\) is given by the nonzero subspaces
of \(T=\Span\{e_1+e_2,\ e_3+e_4,\ e_5+e_6\}\), which is a totally
isotropic \(3\)-space.
\end{corollary}
\section{The line graph in dimension six}
\label{sec:line}
\begin{lemma}\label{lem:weighted}
An independent set in $\Gamma_6$ containing $t$ even-weight vectors and $o$
odd-weight vectors satisfies $3t+4o\leq19$. Its cardinality is at most six;
cardinality six requires $(t,o)=(5,1)$.
\end{lemma}
\begin{proof}
Write $j=(1,1,1,1,1,1)$. The even-weight vectors lie in $j^\perp$, whose
restricted form has radical $\Span(j)$ and rank four. If $o=0$, the Gram
matrix of the selected vectors is $J_t+I_t$, of rank $t$ for even $t$ and
$t-1$ for odd $t$. Thus $t\leq5$.

If $o>0$, fix an odd selected vector $a$. For each even selected vector $e$,
the vector $z_e=e+a$ belongs to $a^\perp$. These vectors are orthonormal and
span a nondegenerate $t$-space $A$. The ambient space $a^\perp$ is
nondegenerate of dimension five. The span $D$ of the vectors $u+a$, over odd
selected vectors $u$, is totally isotropic and orthogonal to $A$ inside
$a^\perp$. Hence $2\dim D\leq5-t$. The odd selected vectors are distinct
members of $a+D$, so $o\leq2^{\dim D}$. Consequently $o\leq4$ for $t=0,1$,
$o\leq2$ for $t=2,3$, and $o\leq1$ for $t=4,5$. These bounds prove both claims.
\end{proof}

\begin{theorem}\label{thm:lines}
$12\leq\chi(\Gamma_6)\leq13$.
\end{theorem}
\begin{proof}
There are $31$ nonzero even-weight and $32$ odd-weight vectors. Their total
weight, using weights three and four respectively, is $221$. By
Lemma~\ref{lem:weighted}, every color class has weight at most $19$, giving
$\chi(\Gamma_6)\geq\lceil221/19\rceil=12$.
For the upper bound, interpret an integer $x$ as the vector whose coordinate
$i$ is its $i$th binary digit. The certificate \texttt{results/line-coloring.json}
partitions the $63$ nonzero vectors into $13$ classes. The independent checker
verifies that every pair of distinct vectors in each class has dot product one.
\end{proof}

For completeness, the thirteen color classes, in the same binary integer
coordinates, are displayed below. They partition $\{1,\ldots,63\}$, and
their within-class dot products can be checked directly.
\[
\begin{gathered}
\{3,29,45,53,57,62\},\quad \{12,20,39,56,59\},\quad
\{15,18,43,52,61\},\\
\{23,24,44,47,48\},\quad \{16,50,51,54,58\},\quad
\{32,46,55,60\},\\
\{13,22,42,49,63\},\quad \{5,19,28,38,41\},\quad
\{8,10,25,40\},\\
\{2,26,30,35\},\quad \{1,7,9,17,33\},\quad
\{6,11,27,34,37\},\quad \{4,14,21,31,36\}.
\end{gathered}
\]

Any hypothetical $12$-coloring has a class of size six. By the lemma it
consists of an odd vector $a$ and five even vectors $e$. The vectors $a$ and
$e+a$ form an orthonormal basis. An isometry therefore maps this class to
$\{1,3,5,9,17,33\}$. Fixing this class reduces the remaining question to an
$11$-coloring of the induced graph on the other $57$ vectors.

\section{The full chromatic number in dimension six}
\label{sec:chromatic}

\begin{theorem}\label{thm:chromatic}
\(\chi(\OG_6^*)=15\).
\end{theorem}

\begin{proof}
Let \(T=\Span(3,12,48)\) in binary integer coordinates, with bit
\(i-1\) representing the coefficient of \(e_i\).  Thus its generators
are \(e_1+e_2,\ e_3+e_4,\ e_5+e_6\).  The form vanishes on \(T\), so
its seven lines, seven planes and \(T\) itself form a \(15\)-clique.
Hence \(\chi(\OG_6^*)\ge15\).

For the upper bound, the file
\path{results/full-15-coloring.json} gives a basis and a color in
\(\{0,\ldots,14\}\) for every nonzero subspace.  The independent
checker reconstructs each span and verifies basis independence,
distinctness, and the dimension counts
\[
63,\ 651,\ 1395,\ 651,\ 63,\ 1 .
\]
These equal the Gaussian binomial coefficients for \(d=1,\dots,6\),
which proves completeness of the catalogue.  For all \(3\,986\,076\)
pairs of distinct entries the checker tests orthogonality on the
basis vectors.  Bilinearity identifies this test with graph
adjacency.  All \(44\,968\) orthogonality edges have distinct
endpoint colors.  This establishes the upper bound
\(\chi(\OG_6^*)\le15\), and therefore equality.
\end{proof}

\begin{remark}\label{rem:precoloring}
To find the certificate, fix distinct colors on the \(15\) nonzero
subspaces \(W\) of \(T\).  Every \(15\)-coloring uses all colors on
this clique, so relabeling colors makes this precoloring an
equivalent instance.  The forbidden colors for a vertex
\(U\not\subseteq T\) are exactly those \(W\subseteq H(U)=U^\perp\cap
T\).  If \(H(U)=T\), then \(U\subseteq T^\perp=T\), a contradiction.
Thus \(\dim H(U)\in\{0,1,2\}\), giving lists of sizes
\[
15-N(\dim H(U))\in\{15,14,11\} .
\]
\end{remark}

\begin{remark}\label{rem:degree-deletion}
A vertex whose residual degree is smaller than its list size can be
deleted and restored greedily after the remaining graph is colored.
Repeated deletion removes \(729\) uncolored vertices and leaves
\(2\,080\) vertices with \(39\,400\) edges.  The list-coloring
encoding has \(28\,600\) variables and \(627\,510\) clauses.  There
is one Boolean variable for each allowed vertex--color pair.  Each
vertex has an at-least-one clause and pairwise at-most-one clauses;
each edge forbids simultaneous use of any color common to its
endpoints' lists.  Hence satisfying assignments are exactly the
proper list colorings of the residual graph.  A satisfying
assignment was extended in reverse deletion order to produce the
full certificate checked above.
\end{remark}

\section{Geometry of the precoloring}
\label{sec:geometry}

In this section we describe the \(15\)-coloring of \(\OG_6^*\)
constructed in Section~\ref{sec:chromatic} in geometric terms.  Keep
the notation of that section: \(T\) is the totally isotropic
\(3\)-space
\[
T=\Span\{e_1+e_2,\ e_3+e_4,\ e_5+e_6\},
\]
\(\mathcal C\) is the \(15\)-clique formed by the nonzero subspaces of
\(T\), and each \(W\in\mathcal C\) carries its own color \(c_W\).

For a vertex \(U\) of \(\OG_6^*\) outside \(\mathcal C\), define
\[
H(U)=U^\perp\cap T,\qquad K(U)=U\cap T.
\]
Since \(T^\perp=T\) and \(U\not\subseteq T\), we have \(H(U)\ne T\), so
\(\dim H(U)\in\{0,1,2\}\).  The following observation is the key to
the precoloring.

\begin{lemma}\label{lem:H-adjacency}
Let \(U\) be a vertex of \(\OG_6^*\) outside \(\mathcal C\) and let
\(W\in\mathcal C\).  Then \(U\) is adjacent to \(W\) if and only if
\(W\subseteq H(U)\).  Consequently, the allowed color list of \(U\)
is
\[
L(U)=\{c_W : 0\ne W\le T,\ W\not\subseteq H(U)\},
\]
and its size is
\[
|L(U)|=\begin{cases}
15,&\dim H(U)=0,\\
14,&\dim H(U)=1,\\
11,&\dim H(U)=2.
\end{cases}
\]
\end{lemma}

\begin{proof}
By definition, \(U\) is adjacent to \(W\) if and only if
\(W\subseteq U^\perp\).  Since \(W\subseteq T\), this is equivalent to
\(W\subseteq U^\perp\cap T=H(U)\).  The list size is
\(15-N(\dim H(U))\), where \(N(0)=0\), \(N(1)=1\), and \(N(2)=4\);
this gives \(15\), \(14\), and \(11\) in the three cases.
\end{proof}

The distribution of the \(2809=2824-15\) vertices outside
\(\mathcal C\), grouped by the triple
\((\dim U,\dim H(U),\dim K(U))\), is recorded in
Table~\ref{tab:profiles}.

\begin{table}[htbp]
\centering
\begin{tabular}{c|r|r|r}
\toprule
\(\dim U\) & \(\dim H(U)\) & \(\dim K(U)\) & Vertices \\
\midrule
\(1\) & \(2\) & \(0\) & \(56\) \\
\(2\) & \(1\) & \(0\) & \(448\) \\
\(2\) & \(2\) & \(1\) & \(196\) \\
\(3\) & \(0\) & \(0\) & \(512\) \\
\(3\) & \(1\) & \(1\) & \(784\) \\
\(3\) & \(2\) & \(2\) & \(98\) \\
\(4\) & \(0\) & \(1\) & \(448\) \\
\(4\) & \(1\) & \(2\) & \(196\) \\
\(4\) & \(2\) & \(3\) & \(7\) \\
\(5\) & \(0\) & \(2\) & \(56\) \\
\(5\) & \(1\) & \(3\) & \(7\) \\
\(6\) & \(0\) & \(3\) & \(1\) \\
\bottomrule
\end{tabular}
\caption{Distribution of the \(2809\) vertices outside the
\(15\)-clique, according to the triple
\((\dim U,\dim H(U),\dim K(U))\).}
\label{tab:profiles}
\end{table}

The \(357\) vertices with \(\dim H(U)=2\) are the most constrained,
having only \(11\) allowed colors.  A direct analysis of the
incidence structure between these vertices and the \(15\) clique
colors shows that the precoloring extends; this is the content of
the certificate described in Section~\ref{sec:chromatic}.

\subsection{Towards a structural description}

The invariants \((\dim U,\dim H(U),\dim K(U))\) do not by themselves
determine the color in the certificate.  Grouping the archived
coloring by this triple produces \(240\) distinct profiles outside
the clique; among these, \(139\) receive more than one color.  The
remaining \(101\) profiles receive a single color.  This is the
precise sense in which the present coloring depends on finer data
than the basic geometric invariants.

Two directions appear promising for a structural description.

\begin{enumerate}
\item \emph{Enrichment of the invariants.}  The profiles that receive
a single color suggest that the triple
\((\dim U,\dim H(U),\dim K(U))\) captures part of the coloring rule.
One may look for an additional invariant, such as the position of
\(K(U)\) inside \(H(U)\) or the structure of the quotient
\(U/(U\cap H(U))\), that separates the multi-color profiles.

\item \emph{The stabilizer of \(T\).}  The orthogonal group of
\(V_6\) acts on \(\OG_6^*\) and permutes the \(15\) clique colors.
A coloring invariant under the stabilizer of \(T\) in this group
would reduce the problem to a finite group-theoretic computation.
This direction has not been pursued in the present certificate.
\end{enumerate}

A complete structural proof that \(\chi(\OG_6^*)=15\), avoiding the
finite computation, remains open.  We record this as
Problem~\ref{prob:structural} in Section~\ref{sec:open}.

\section{Higher dimensions and open problems}
\label{sec:higher}
The clique formula of Section~\ref{sec:clique} holds in all dimensions,
while the exact chromatic results of Sections~\ref{sec:line} and
\ref{sec:chromatic} concern dimension six. This section records a
reduction for higher-dimensional coloring problems. We first give a degree-versus-list deletion
rule that reduces a \(k\)-coloring problem in dimension \(n\) to a
smaller residual instance, without changing \(k\)-colorability.
We then apply it to \(n=7\), where a \(16\)-clique is available and
the rule removes all vertices of dimension at least four, leaving a
residual instance of \(14\,589\) vertices and \(953\,057\) edges.
Finally, we state the open problems that remain: the exact
chromatic number for \(n\ge7\), the question of whether
\(\chi>\omega\) holds in infinitely many dimensions, and the structural
question of whether an optimal \(15\)-coloring in dimension six
can be described without finite computation.
\begin{proposition}\label{prop:dimension-deletion}
Let \(K\) be a fixed \(k\)-clique in \(\OG_n^*\).  If \(N(n-d)<k\),
every uncolored \(d\)-dimensional vertex can be deleted by the
degree-versus-list rule.  Consequently, \(k\)-colorability of the
residual instance after such deletions is equivalent to
\(k\)-colorability of the full graph.
\end{proposition}
\begin{proof}
For a \(d\)-space \(U\), the orthogonal complement has dimension
\(n-d\).  Every neighbor of \(U\) is a nonzero subspace of
\(U^\perp\), so \(\deg U\leq N(n-d)\).  If \(f\) fixed clique
vertices are adjacent to \(U\), then its residual degree is at most
\(N(n-d)-f\), while its list size is \(k-f\).  The stated inequality
therefore guarantees an available color on restoring \(U\).  Apply
this argument in reverse deletion order.
\end{proof}

For \(n=7\), the nonzero subspaces of \(T=\Span(3,12,48)\) together
with \(\Span(64)\) give a \(16\)-clique.  Since \(N(3)=15<16\),
Proposition~\ref{prop:dimension-deletion} removes all \(14\,606\)
vertices of dimension at least four from the \(29\,211\)-vertex
graph.  The remaining uncolored instance has \(14\,589\) vertices
and \(953\,057\) edges.  The direct pairwise list-color encoding
would have \(214\,355\) variables and \(12\,984\,679\) clauses.
These exact instance sizes guide further geometric reduction and
resource-bounded search.  They are measured preprocessing counts;
no \(16\)-coloring or checked nonexistence certificate for
dimension seven is supplied.

Problem~4.3 remains open for \(n\geq7\), as does Problem~4.4 of
\cite{Nikandish2026}.  Within dimension six, a structural description
of an optimal full-subspace coloring remains an immediate question.

The negative results of Section~\ref{sec:geometry} --- in particular
Proposition~\ref{prop:lift-features} and
Proposition~\ref{prop:stabT-orbit} --- rule out colorings depending
only on the exact profiles or the specified lift features, as well as
colorings invariant under the stabilizer of \(T\).
A structural proof of
\(\chi(\OG_6^*)=15\) must therefore use auxiliary data that break
this symmetry and retain enough placement information to separate
adjacent vertices.

\begin{problem}\label{prob:structural}
Can one specify auxiliary geometric data and a nontrivial proper subgroup
\(H<\operatorname{Stab}_{O_6(\F)}(T)\), and construct an explicit proper
\(15\)-coloring \(c\) from those data such that
\(c(hU)=h\cdot c(U)\) for \(h\in H\), where the color labels are the
nonzero subspaces of \(T\) with their induced \(H\)-action?

The construction should prove availability and separation without
looking up the archived coloring certificate. Such a construction would
give a conceptual proof of \(\chi(\OG_6^*)=15\).
The checked \(12\)-coloring of \(\Gamma_6\) establishes its chromatic number;
equivariance under the full orthogonal group is not established here.
\end{problem}

\begin{remark}\label{prob:gamma-six}
The line-graph question is settled by Theorem~\ref{thm:lines}:
\(\chi(\Gamma_6)=12\), whereas \(\chi(\OG_6^*)=15\).
Thus passing from lines to all nonzero subspaces increases the
chromatic number by exactly three in dimension six.
\end{remark}

\section{Computational certificates and formal verification}\label{sec:verification}
The positive coloring certificates are checked using integer arithmetic and
the coordinate definition of orthogonality. The full checker imports neither
the search program nor a SAT solver. It establishes enumeration completeness
from distinctness and Gaussian counts, then checks every pair of vertices.
Three deliberately invalid certificates, respectively omitting a vertex,
duplicating a subspace, and identifying the colors of adjacent coordinate
lines, are rejected. The line-coloring checker independently verifies its
partition and all within-class pairs.

The radical/quotient clique theorem has an existing Lean formalization.
Its source is \path{formal/NikandishClique.lean}, with theorem
\texttt{NikandishClique.result} in its documented namespace. The exact
upstream revision, Lean version and dependencies are recorded in
\texttt{formal/UPSTREAM.json}. The new chromatic certificate is currently
verified by the independent integer checker; its Lean verification is being
developed separately from the archived clique theorem.

The public package~\cite{Package} contains the literal witnesses, independent
checkers, search records, and SHA-256 bindings. The new geometric profiler
checks every assigned color against its $H(U)$-determined list. The
dimension-seven reduction is also replayed in dimension six, where it
reproduces the established core size and Boolean encoding counts.

\begin{thebibliography}{9}
\bibitem{Nikandish}
R. Nikandish, \emph{Annihilating-Ideal Graphs and Orthogonality Graphs over
$\mathbb F_2$}, arXiv:2609.22769v1, 2026.
\bibitem{Package}
H. Ma, R. Nikandish and W. Zhang, \emph{Binary subspace orthogonality graphs:
research and reproducibility workspace},
\url{https://github.com/the-omega-institute/binary-subspace-orthogonality}.
\end{thebibliography}
\end{document}
